\ticket{16}{Нейронные сети прямого распространения. Архитектуры MLP и RBF: структура сетей, виды функций активации, алгоритмы обучения, проблема переобучения и локальных минимумов}

\begin{examplan}
    \item Определить сети прямого распространения.
    \item Описать архитектуру MLP и функции активации.
    \item Объяснить обучение через backpropagation.
    \item Описать RBF-сети и их обучение.
    \item Рассказать о переобучении и локальных минимумах.
\end{examplan}

\subsection*{Короткая версия для письменного ответа}

\textbf{Нейронные сети прямого распространения} --- сети, в которых информация движется от входного слоя к выходному без циклов и обратных связей. Они используются для классификации, регрессии и аппроксимации функций. Основные архитектуры в билете: MLP и RBF.

\subsection*{MLP}

\textbf{Многослойный перцептрон} состоит из входного слоя, одного или нескольких скрытых слоев и выходного слоя. В полносвязном слое каждый нейрон получает выходы предыдущего слоя.

Нейрон вычисляет:
\[
z=\sum_i w_i x_i+b,\qquad a=\varphi(z),
\]
где $w_i$ --- веса, $b$ --- смещение, $\varphi$ --- функция активации.

Функции активации:
\begin{itemize}
    \item sigmoid --- часто для бинарной классификации, но страдает затухающими градиентами;
    \item tanh --- центрирована около нуля, но тоже может насыщаться;
    \item ReLU --- популярна в скрытых слоях, проста и уменьшает затухание градиента;
    \item Leaky ReLU --- снижает проблему "мертвых" ReLU;
    \item linear --- часто для регрессии;
    \item softmax --- для многоклассовой классификации.
\end{itemize}

Без нелинейных функций активации многослойная сеть сводилась бы к линейной модели.

\subsection*{Обучение MLP}

Обучение состоит в настройке весов и смещений для минимизации функции потерь.

Основной алгоритм --- \textbf{обратное распространение ошибки}:
\begin{enumerate}
    \item случайно инициализировать веса;
    \item выполнить прямой проход и получить прогноз;
    \item вычислить ошибку;
    \item распространить градиенты от выхода к входу по правилу цепочки;
    \item обновить веса, например $w\leftarrow w-\eta \frac{\partial L}{\partial w}$;
    \item повторять по эпохам.
\end{enumerate}

Популярные оптимизаторы: SGD, mini-batch SGD, RMSprop, Adam.

\subsection*{RBF-сети}

\textbf{RBF-сеть} обычно имеет три слоя: входной, скрытый слой радиальных базисных функций и линейный выходной слой.

Скрытый нейрон имеет центр $c_i$ и ширину $\sigma_i$. Часто используется гауссова функция:
\[
\varphi_i(x)=\exp\left(-\frac{\|x-c_i\|^2}{2\sigma_i^2}\right).
\]

Активация максимальна, когда вход близок к центру, поэтому RBF выполняет локальную аппроксимацию.

Обучение RBF часто двухэтапное:
\begin{enumerate}
    \item определить центры, например алгоритмом $k$-means;
    \item выбрать ширины;
    \item обучить веса выходного линейного слоя как задачу линейной регрессии или градиентным методом.
\end{enumerate}

\subsection*{MLP и RBF}

\begin{description}
    \item[MLP] обычно выполняет глобальную аппроксимацию, обучается backpropagation, может быть глубоким.
    \item[RBF] выполняет локальную аппроксимацию вокруг центров, часто обучается быстрее, но может требовать больше скрытых нейронов.
\end{description}

\subsection*{Переобучение}

\textbf{Переобучение} возникает, когда сеть хорошо запоминает обучающую выборку, но плохо обобщает на новые данные. Причины: слишком сложная модель, мало данных, слишком долгое обучение, шум в данных.

Методы борьбы:
\begin{itemize}
    \item L1/L2-регуляризация;
    \item dropout;
    \item ранняя остановка;
    \item аугментация данных;
    \item уменьшение сложности модели;
    \item валидационная выборка и кросс-валидация.
\end{itemize}

\subsection*{Локальные минимумы}

Функция потерь нейронной сети невыпукла. Градиентные методы могут попадать в локальные минимумы или долго проходить седловые точки. На практике помогают удачная инициализация, mini-batch шум, оптимизаторы с моментом, Adam/RMSprop и несколько запусков обучения.

\begin{important}
    \item Feedforward-сеть не имеет циклов: сигнал идет от входа к выходу.
    \item MLP использует линейные преобразования и нелинейные активации.
    \item Backpropagation вычисляет градиенты по правилу цепочки.
    \item RBF использует расстояния до центров и радиальные функции.
    \item Переобучение контролируют регуляризацией, dropout и ранней остановкой.
\end{important}
